\chapter{Salpingectomy}

\section*{Introduction \& Clinical Indications}
Salpingectomy is a surgical procedure involving the complete excision of one or both fallopian tubes, which are essential structures in the female reproductive system responsible for transporting ova from the ovaries to the uterus. This intervention is clinically significant in various scenarios, including the management of ectopic pregnancies, where the fertilized egg implants outside the uterus, often within a fallopian tube. Salpingectomy is also indicated for severe tubal infections, such as salpingitis or hydrosalpinx, which can lead to chronic pain and infertility. Additionally, this procedure serves as a permanent contraceptive method, effectively preventing future pregnancies. In recent years, salpingectomy has gained prominence as a prophylactic measure to reduce the risk of ovarian cancer, particularly in women with genetic predispositions such as BRCA1 or BRCA2 mutations. The procedure can be performed unilaterally or bilaterally, depending on the clinical indication and the patient's reproductive goals, making it a vital component of women's health management.

\begin{itemize}
  \item Fallopian tube removal
  \item Tubectomy
  \item Unilateral tubal excision
  \item Bilateral salpingectomy (BS)
  \item Opportunistic salpingectomy (for cancer risk reduction)
  \item Risk-reducing salpingectomy (RRS)
\end{itemize}

The primary surgical principle of salpingectomy is the complete excision of the affected fallopian tube(s) to achieve specific therapeutic, prophylactic, or contraceptive goals. In cases of ectopic pregnancy, the main objective is to remove the non-viable gestation and achieve hemostasis, thereby preventing life-threatening tubal rupture and preserving maternal health. For severe tubal infections, such as pyosalpinx or hydrosalpinx, the removal of the inflamed or fluid-filled tube alleviates chronic pain and can improve fertility outcomes by eliminating a damaged, non-functional tube that may impair in vitro fertilization (IVF) success.

As a method of permanent contraception, bilateral salpingectomy interrupts the pathway for sperm to reach the ovum and for the ovum to reach the uterus, providing a highly effective and irreversible sterilization method. Furthermore, the procedure has gained significant importance in cancer risk reduction, particularly for high-grade serous ovarian cancer. Emerging evidence suggests that a substantial proportion of these aggressive cancers originate in the fimbriated end of the fallopian tube, specifically from serous tubal intraepithelial carcinoma (STIC) lesions. By excising the entire fallopian tube, especially the fimbriae, the presumed site of origin for these cancers is eliminated, offering a substantial reduction in risk for women, particularly those with BRCA1 or BRCA2 mutations, and increasingly, as an opportunistic measure during other benign pelvic surgeries. The rationale is to maximize cancer risk reduction by removing the entire tube, including the fimbria.

The history of salpingectomy is closely linked to the evolution of gynecologic surgery and the understanding of female reproductive pathology. Early attempts at removing diseased fallopian tubes date back to the late 19th century, driven by the urgent need to address life-threatening conditions such as ruptured ectopic pregnancies and severe pelvic inflammatory disease (PID). Robert Lawson Tait, a pioneering Scottish surgeon, is credited with performing one of the first successful abdominal surgeries for a ruptured ectopic pregnancy in 1883, often involving salpingectomy. These initial procedures were conducted via open laparotomy, which, while life-saving, carried significant morbidity and prolonged recovery times.

The latter half of the 20th century witnessed a revolutionary shift with the advent of minimally invasive surgery. Kurt Semm, a German gynecologist, made pivotal contributions in the 1970s and 1980s by developing and popularizing laparoscopic techniques for a wide array of gynecologic procedures, including salpingectomy. This innovation dramatically reduced recovery times, postoperative pain, and surgical morbidity, making salpingectomy a much safer and more accessible procedure. More recently, a paradigm shift in gynecologic oncology occurred with the growing understanding that many "ovarian" cancers, specifically high-grade serous carcinomas, originate in the fimbriated end of the fallopian tube rather than the ovary itself. This led to the concept of "opportunistic salpingectomy," where the fallopian tubes are removed during other pelvic surgeries (e.g., hysterectomy for benign indications) primarily to reduce the risk of ovarian cancer. This practice, championed by organizations like the Society of Gynecologic Oncology, has gained widespread acceptance in the 21st century, marking a significant evolution in the prophylactic role of salpingectomy.

Salpingectomy is indicated for a range of specific clinical conditions, each addressing a distinct medical necessity:

\begin{itemize}
  \item **Ectopic Pregnancy:** To remove a pregnancy implanted outside the uterus, most commonly within the fallopian tube, especially if it is ruptured, symptomatic, or unresponsive to medical management with methotrexate.
  \item **Tubal Infection (Salpingitis/Pyosalpinx):** To treat severe, recurrent, or chronic infections of the fallopian tube, often leading to abscess formation (pyosalpinx) or fluid accumulation (hydrosalpinx), which can cause intractable pain and infertility.
  \item **Hydrosalpinx:** To remove a distally occluded, fluid-filled fallopian tube, particularly in the context of infertility treatment, as the fluid can be embryotoxic and significantly impair in vitro fertilization (IVF) success rates.
  \item **Tubal Torsion:** To remove a fallopian tube that has twisted on its vascular pedicle, leading to ischemia, necrosis, and acute, severe pelvic pain, often requiring emergent intervention.
  \item **Permanent Contraception (Sterilization):** As an elective procedure for women who desire highly effective and irreversible birth control, offering a definitive method of preventing future pregnancies.
  \item **Risk-Reducing Surgery for Ovarian Cancer:** For women with a high genetic predisposition (e.g., BRCA1/2 mutations, Lynch syndrome) to significantly reduce the risk of developing high-grade serous ovarian cancer.
  \item **Opportunistic Salpingectomy:** Performed incidentally during other benign pelvic surgeries (e.g., hysterectomy, myomectomy) in the general population to reduce the lifetime risk of ovarian cancer.
  \item **Benign Tubal Masses or Cysts:** To remove non-cancerous growths or cysts originating from the fallopian tube that are causing symptoms, are rapidly growing, or are of uncertain etiology.
  \item **Endometriosis:** In severe cases where endometriosis extensively involves the fallopian tube, causing significant pain or tubal damage, salpingectomy may be performed as part of a broader surgical management strategy.
\end{itemize}

Salpingectomy holds a versatile clinical positioning, serving as both a primary and a secondary procedure, depending on the specific clinical context and urgency. It is frequently a primary, definitive surgical treatment for acute, life-threatening conditions such as a ruptured ectopic pregnancy, where immediate intervention is paramount to control hemorrhage and stabilize the patient. Similarly, for severe tubal infections (e.g., pyosalpinx) unresponsive to antibiotic therapy, or for acute tubal torsion, salpingectomy is typically the primary and most effective therapeutic approach.

In other scenarios, salpingectomy may be considered a secondary or elective procedure. For unruptured ectopic pregnancies, medical management with methotrexate is often the initial primary approach, with salpingectomy reserved for cases of treatment failure, contraindications to medical therapy, or patient preference for surgical resolution. As a method of permanent contraception or for cancer risk reduction, salpingectomy is an elective, primary procedure chosen by the patient after thorough counseling. When performed opportunistically during another primary surgery (e.g., hysterectomy for uterine fibroids or abnormal uterine bleeding), the salpingectomy itself is secondary to the main indication for the initial surgery, but it serves a primary prophylactic role for ovarian cancer risk reduction, adding significant long-term health benefits. This dual role underscores its importance in gynecologic surgery.

\section*{Relevant Surgical Anatomy}
The fallopian tubes, also referred to as oviducts or uterine tubes, are paired muscular structures approximately 10-12 cm in length, extending bilaterally from the uterus to the ovaries. Each tube is anatomically divided into four segments: the infundibulum (with fimbriae), ampulla, isthmus, and intramural (interstitial) portion. The fimbriated end of the infundibulum is positioned near the ovary, facilitating the capture of the ovum post-ovulation. The tubes are supported by the mesosalpinx, a peritoneal fold of the broad ligament, which contains their neurovascular supply.

The arterial supply to the fallopian tubes arises from anastomoses between the ovarian artery (a direct branch of the abdominal aorta) and the tubal branch of the uterine artery (a branch of the internal iliac artery). Venous drainage parallels the arterial supply, with veins draining into the ovarian veins (right to the inferior vena cava, left to the left renal vein) and the uterine venous plexus (to internal iliac veins). Lymphatic drainage primarily follows the ovarian vessels to the para-aortic lymph nodes, with some drainage also directed to the internal iliac nodes. Innervation is autonomic, originating from the ovarian plexus and the uterovaginal plexus. Key surgical landmarks include the round ligament anteriorly, the infundibulopelvic ligament (containing the ovarian vessels) laterally, and the utero-ovarian ligament medially, which assist in delineating the adnexal structures during pelvic surgery. Understanding these anatomical relationships is crucial for the safe and complete excision of the fallopian tubes while preserving adjacent structures.

\section*{Preoperative Workup \& Preparation}
Thorough preoperative workup and preparation are paramount to ensure patient safety and optimize surgical outcomes for salpingectomy. This comprehensive process typically includes:

\begin{itemize}
  \item **Comprehensive Medical Evaluation:** A detailed medical history and physical examination are performed to assess the patient's overall health status, identify any pre-existing comorbidities (e.g., cardiac, pulmonary, renal disease), and evaluate the specific indication for surgery. This includes a review of past surgical history, allergies, and current medications.
  \item **Laboratory Investigations:** Routine blood tests are ordered, including a complete blood count (CBC) to assess for anemia or infection, blood type and screen for potential transfusion needs, a coagulation profile (PT/INR, PTT) to evaluate bleeding risk, and renal and liver function tests. For all women of reproductive age, a quantitative serum beta-human chorionic gonadotropin (beta-hCG) test is mandatory to definitively rule out an intrauterine pregnancy.
  \item **Imaging Studies:** Transvaginal ultrasound is frequently performed to confirm the diagnosis (e.g., ectopic pregnancy, hydrosalpinx, tubal torsion) and to meticulously assess pelvic anatomy, including the uterus and ovaries. Other imaging modalities, such as computed tomography (CT) or magnetic resonance imaging (MRI), may be utilized for complex cases, suspected malignancy, or to delineate extensive adhesions.
  \item **Anesthesia Consultation:** A dedicated consultation with an anesthesiologist is conducted to evaluate the patient's fitness for general anesthesia, discuss potential anesthetic risks, and formulate a tailored anesthesia plan.
  \item **Medication Management:** All current medications are meticulously reviewed. Anticoagulants (e.g., warfarin, DOACs) or antiplatelet agents (e.g., aspirin, clopidogrel) may need to be temporarily discontinued or adjusted according to established protocols to minimize bleeding risk.
  \item **NPO Status:** Patients are strictly instructed to be NPO (nil per os - nothing by mouth) for a specified period (typically 6-8 hours for solids, 2 hours for clear liquids) before surgery to significantly reduce the risk of pulmonary aspiration during anesthesia induction.
  \item **Bowel Preparation:** In cases where extensive pelvic adhesions, endometriosis, or a high likelihood of bowel involvement is anticipated, a mild mechanical bowel preparation may be recommended to reduce bowel contents and facilitate safer surgical access.
  \item **Prophylactic Antibiotics:** Intravenous broad-spectrum antibiotics (e.g., a first-generation cephalosporin like cefazolin) are typically administered within 60 minutes prior to the surgical incision to reduce the risk of surgical site infection.
  \item **Informed Consent:** A comprehensive discussion with the patient is undertaken, detailing the procedure, its specific indications, potential benefits, inherent risks (including general surgical and procedure-specific complications), available alternatives, and the long-term implications (e.g., fertility, contraception, cancer risk reduction). Fully informed consent is then formally obtained and documented.
\end{itemize}

While salpingectomy is a common and generally safe procedure, certain conditions may contraindicate its performance or necessitate significant precautions.

\textbf{Situations requiring treatment, optimization, or an alternative plan:}
\begin{itemize}
  \item **Uncontrolled Coagulopathy:** Severe, uncorrected bleeding disorders pose an unacceptably high risk of intraoperative and postoperative hemorrhage, which cannot be safely managed.
  \item **Severe Cardiopulmonary Instability:** Patients with unstable cardiac conditions (e.g., recent myocardial infarction, uncontrolled arrhythmias) or severe respiratory compromise (e.g., acute respiratory distress syndrome) that preclude safe administration of general anesthesia or tolerance of pneumoperitoneum.
  \item **Patient Refusal:** The patient's informed and autonomous decision to decline the procedure, after comprehensive counseling, is an absolute contraindication.
\end{itemize}

\textbf{Relative Contraindications \& Precautions:}
\begin{itemize}
  \item **Extensive Intra-abdominal Adhesions:** Previous abdominal surgeries, severe pelvic inflammatory disease, or endometriosis can lead to dense adhesions, making laparoscopic access challenging, increasing operative time, and significantly elevating the risk of iatrogenic injury to adjacent organs (e.g., bowel, bladder). An open approach may be preferred or necessary.
  \item **Morbid Obesity:** While not an absolute contraindication, severe obesity can make laparoscopic access and visualization technically more demanding, potentially increasing operative time, anesthetic risks, and the risk of port-site complications.
  \item **Active Systemic Infection or Sepsis:** Although salpingectomy may be indicated for localized pelvic infection (e.g., pyosalpinx), active systemic sepsis requires aggressive medical stabilization before elective or semi-elective surgery.
  \item **Desire for Future Fertility (for bilateral salpingectomy):** Bilateral salpingectomy causes permanent loss of tubal natural conception and requires explicit fertility counseling; future pregnancy may require IVF. Alternatives should be discussed when clinically appropriate.
  \item **Large Uterine Fibroids or Pelvic Masses:** These can significantly distort normal pelvic anatomy, impede surgical access, and obscure visualization of the fallopian tubes, requiring careful preoperative planning and potentially an alternative surgical approach.
  \item **Intrauterine Pregnancy (other than ectopic):** If a viable intrauterine pregnancy is present, salpingectomy is generally contraindicated unless there is a separate, compelling, and life-threatening indication for the procedure.
  \item **Severe Endometriosis:** While salpingectomy can be part of endometriosis management, severe, diffuse endometriosis can create a "frozen pelvis," making dissection extremely difficult and risky.
\end{itemize}

\section*{Surgical Technique \& Procedure}
Salpingectomy is most commonly performed using a minimally invasive laparoscopic approach, which offers reduced recovery times and morbidity compared to open laparotomy. However, an open approach may be necessary in complex cases, emergencies with massive hemorrhage, or when extensive adhesions are present.

\begin{itemize}
  \item **Anesthesia and Patient Positioning:** The patient is placed under general endotracheal anesthesia. She is then positioned in the dorsal lithotomy position, often with a steep Trendelenburg tilt, to allow gravity to displace the bowel superiorly, thereby optimizing visualization of the pelvic organs. The abdomen and perineum are prepped and draped in a sterile fashion.
  \item **Incisions and Trocar Placement:** A small infraumbilical incision (typically 10-12 mm) is made for the primary trocar, through which the laparoscope is inserted. The abdomen is then insufflated with carbon dioxide to create a pneumoperitoneum, establishing a working space. Two or three additional smaller incisions (5 mm) are made in the lower abdomen (e.g., suprapubic, bilateral lower quadrants) for accessory trocars, facilitating the introduction of surgical instruments.
  \item **Pelvic Exploration and Identification:** The surgeon systematically explores the pelvic and abdominal cavities to identify the uterus, ovaries, and fallopian tubes. The affected fallopian tube is carefully identified, and its relationship to surrounding structures, including the ovary and broad ligament, is assessed.
  \item **Dissection and Hemostasis of the Mesosalpinx:** Using laparoscopic graspers and advanced energy devices (e.g., bipolar cautery, harmonic scalpel, or ligating clips), the mesosalpinx (the peritoneal fold supporting the fallopian tube) is carefully coagulated and transected along its entire length. This step meticulously separates the fallopian tube from its vascular supply and the broad ligament, ensuring adequate hemostasis.
  \item **Tubal Transection:** The fimbriated (distal) end of the fallopian tube is sealed and divided. Subsequently, the cornual (proximal) end of the tube, where it enters the uterus, is also sealed and transected, ensuring complete removal of the entire tube, including the interstitial portion if indicated for cancer risk reduction.
  \item **Specimen Retrieval:** Once completely detached and free, the fallopian tube is carefully placed into an endoscopic retrieval bag. The bag is then extracted through one of the larger trocar sites, often the umbilical port, to prevent contamination of the abdominal cavity.
  \item **Inspection, Irrigation, and Hemostasis:** The surgical field is meticulously inspected for any bleeding points, which are then coagulated. The pelvis is thoroughly irrigated with sterile saline to remove any blood clots or debris, and the fluid is aspirated.
  \item **Closure:** The trocars are removed under direct visualization. The fascial defects at the larger port sites (typically >10 mm) are closed with absorbable sutures to prevent incisional hernia formation. All skin incisions are then closed with subcuticular sutures, surgical adhesive, or sterile strips. Drains are typically not required for uncomplicated salpingectomy.
\end{itemize}

In rare instances, such as extensive adhesions, massive hemorrhage, or suspected malignancy requiring wider exposure, an open laparotomy may be performed, involving a larger abdominal incision (e.g., Pfannenstiel or vertical midline).

\section*{Complications \& Risks}
As with any surgical procedure, salpingectomy carries potential risks, which can be broadly categorized into general surgical complications and those specific to the procedure.

\textbf{General Surgical Complications:}
\begin{itemize}
  \item **Anesthesia Complications:** These include adverse reactions to anesthetic agents, postoperative nausea and vomiting, aspiration pneumonia, respiratory depression, and, rarely, severe cardiac events (e.g., myocardial infarction) or cerebrovascular accidents (stroke).
  \item **Bleeding:** Intraoperative hemorrhage, which may necessitate blood transfusion, or postoperative hematoma formation at the surgical site or within the abdominal cavity.
  \item **Infection:** Surgical site infection (wound infection), pelvic infection (e.g., abscess, cellulitis), or more generalized peritonitis or sepsis.
  \item **Organ Injury:** Accidental damage to adjacent organs during trocar insertion or dissection, including perforation of the bowel, injury to the bladder, ureters, or major blood vessels.
  \item **Nerve Injury:** Damage to peripheral nerves (e.g., femoral nerve, obturator nerve) from improper patient positioning or prolonged pressure, leading to temporary or permanent numbness, weakness, or pain.
  \item **Adhesion Formation:** Development of scar tissue within the abdomen, which can lead to chronic pelvic pain, bowel obstruction, or infertility in the future.
  \item **Deep Vein Thrombosis (DVT) and Pulmonary Embolism (PE):** Formation of blood clots in the deep veins, which can dislodge and travel to the lungs, a potentially life-threatening complication.
\end{itemize}

\textbf{Procedure-Specific Risks:}
\begin{itemize}
  \item **Incomplete Removal of Ectopic Pregnancy:** If a portion of the tubal pregnancy is inadvertently left behind, there is a risk of persistent ectopic pregnancy, necessitating further medical (e.g., methotrexate) or surgical intervention.
  \item **Ovarian Damage:** Injury to the adjacent ovary, particularly if the fallopian tube is densely adherent to it or if a salpingo-oophorectomy (removal of both tube and ovary) is performed, potentially affecting ovarian function or future fertility.
  \item **Chronic Pelvic Pain:** Persistent or new-onset pain in the pelvic region, which can sometimes occur after any pelvic surgery, though it is rare after an uncomplicated salpingectomy.
  \item **Infertility (after bilateral salpingectomy):** Bilateral removal of the fallopian tubes results in permanent and irreversible infertility, as ova can no longer reach the uterus for fertilization.
  \item **Regret (after sterilization):** Some patients may experience psychological distress or regret after permanent sterilization, particularly if life circumstances or reproductive desires change.
  \item **Failure of Sterilization:** Although exceedingly rare (less than 0.1\%), there is a minute chance of pregnancy occurring after bilateral salpingectomy due to recanalization, fistula formation, or incomplete tubal occlusion.
  \item **Incisional Hernia Formation:** Development of a hernia at the port sites, especially at larger trocar sites if fascial closure is inadequate, requiring potential re-operation.
  \item **Persistent Trophoblastic Disease:** In cases of ectopic pregnancy, rarely, residual trophoblastic tissue can persist and require further monitoring and treatment.
\end{itemize}

\section*{Postoperative Care \& Recovery}
The postoperative recovery timeline following salpingectomy is generally rapid, especially with a laparoscopic approach, but varies based on the individual patient, the complexity of the surgery, and the underlying indication.

Immediately after surgery, the patient is transferred to the Post-Anesthesia Care Unit (PACU) for close monitoring of vital signs, pain levels, and recovery from general anesthesia. Pain management is initiated, typically with intravenous analgesics, transitioning to oral medications as tolerated. Nausea and vomiting, common side effects of anesthesia, are managed with antiemetics. Once stable, awake, and pain is controlled, the patient is transferred to a recovery room or inpatient ward. For most laparoscopic salpingectomies, patients are discharged home within 24 hours, often on the same day. Early ambulation is strongly encouraged to prevent complications such as deep vein thrombosis and to aid in the resolution of gas pain caused by residual carbon dioxide from the pneumoperitoneum. A light diet is usually tolerated shortly after surgery, progressing to a regular diet as comfort allows.

Over the first 1-2 weeks, patients typically experience mild to moderate pain at the incision sites and generalized abdominal discomfort, which can be effectively managed with prescribed oral pain medication (e.g., NSAIDs, acetaminophen, or short-term opioids). Shoulder pain, resulting from diaphragmatic irritation by residual carbon dioxide, is common but usually resolves within a few days. Activity restrictions include avoiding heavy lifting (generally over 10-15 pounds), strenuous exercise, and sexual intercourse for 2-4 weeks to allow for adequate internal healing and to prevent wound complications. Incisions should be kept clean and dry. Long-term recovery involves a gradual return to normal daily activities, typically within 2-4 weeks for laparoscopic surgery. For patients undergoing salpingectomy for ectopic pregnancy, serial beta-hCG levels may be monitored until undetectable to ensure complete resolution. For those undergoing sterilization or cancer risk reduction, the long-term impact relates to permanent contraception or a significantly reduced cancer risk, respectively, with no further specific recovery milestones beyond general health maintenance.

Understanding the results of a salpingectomy involves a comprehensive review of the intraoperative surgical findings and the definitive pathology report on the resected tissue. During the procedure, the surgeon meticulously documents the appearance of the fallopian tube(s) and surrounding pelvic structures. For an ectopic pregnancy, this includes noting the size, location, and integrity of the tubal pregnancy, the presence of hemoperitoneum, and the condition of the contralateral tube and ovaries. In cases of infection, the surgeon observes signs of inflammation, adhesions, pus (pyosalpinx), or fluid accumulation (hydrosalpinx). For risk-reducing salpingectomy, the surgeon confirms the complete removal of the entire tube, including the fimbriated end. Any other incidental findings, such as endometriosis, fibroids, or ovarian cysts, are also documented.

The resected fallopian tube(s) are sent to the pathology laboratory for gross and microscopic examination. The pathology report provides the definitive diagnosis. For an ectopic pregnancy, it confirms the presence of chorionic villi within the tubal lumen. For infection, it describes inflammatory changes, possibly with abscess formation or chronic salpingitis. If the procedure was performed for cancer risk reduction, the pathologist meticulously examines the entire tube, with particular attention to the fimbriated end, for any precancerous lesions such as serous tubal intraepithelial carcinoma (STIC) or early invasive carcinoma. For sterilization, the report simply confirms the presence of normal fallopian tube tissue. These pathology results are crucial for guiding subsequent management, such as further monitoring for ectopic pregnancy, additional treatment for infection, or genetic counseling and surveillance recommendations for cancer risk.

A normal and successful postoperative course following salpingectomy is characterized by an uneventful recovery and the achievement of the intended clinical objective. Patients can expect a gradual resolution of pain, which is typically managed effectively with oral analgesics and progressively diminishes over the first few days to weeks. Incision sites should heal cleanly, without signs of infection such as increasing redness, swelling, warmth, or purulent discharge. Any mild bruising or swelling around the port sites is normal and resolves spontaneously. Shoulder pain from residual carbon dioxide usually subsides within 24-48 hours.

For patients who underwent salpingectomy for an ectopic pregnancy, a normal course includes a steady decline in serum beta-hCG levels to undetectable, confirming complete resolution of the pregnancy. For those with infection, symptoms like pain and fever should resolve, indicating successful eradication of the infection. Patients undergoing sterilization or cancer risk reduction will experience no specific disease-related symptoms post-procedure. Most individuals are able to resume light daily activities within a few days and return to full normal activities, including work and exercise, within 2 to 4 weeks, depending on the individual's healing capacity and the surgical approach. A normal course signifies effective treatment or prevention, leading to an improved quality of life or health outcome without significant complications.

Postoperative care and follow-up are essential to monitor recovery, address any concerns, and ensure the long-term success and well-being of the patient after salpingectomy.

\begin{itemize}
  \item **Postoperative Office Visit:** A follow-up appointment with the surgeon is typically scheduled 1-2 weeks after the procedure to assess wound healing, discuss the definitive pathology report, and address any lingering symptoms or questions the patient may have.
  \item **Suture/Staple Removal:** If non-absorbable sutures or surgical staples were used for skin closure, they will be removed at the follow-up visit. Absorbable sutures dissolve on their own.
  \item **Beta-hCG Monitoring:** For patients who underwent salpingectomy for an ectopic pregnancy, serial quantitative beta-hCG levels are often monitored weekly until they are undetectable to confirm complete resolution and rule out persistent trophoblastic disease.
  \item **Contraception Counseling:** For patients who underwent bilateral salpingectomy for sterilization, comprehensive contraception counseling may be provided to discuss other aspects of reproductive health, including sexually transmitted infection prevention.
  \item **Genetic Counseling and Surveillance:** For patients with a genetic predisposition to ovarian cancer who underwent risk-reducing salpingectomy, ongoing genetic counseling and personalized surveillance recommendations will be discussed.
  \item **Activity Resumption Guidance:** Clear instructions are provided on when to safely resume strenuous activities, exercise, and sexual intercourse, typically after 2-4 weeks, based on individual healing and comfort.
  \item **Warning Signs Education:** Patients are thoroughly educated on warning signs that warrant immediate medical attention, such as persistent or worsening abdominal pain, fever (temperature >100.4$^{\circ}$F or 38$^{\circ}$C), heavy vaginal bleeding, foul-smelling vaginal discharge, persistent nausea and vomiting, or signs of wound infection (increasing redness, swelling, pus).
\end{itemize}

Ongoing routine gynecologic care, including regular check-ups and age-appropriate cancer screenings, remains important for all patients.

\section*{Outcomes \& Prognosis}
Salpingectomy is a procedure with a varied epidemiological profile, largely dependent on its specific indication. Ectopic pregnancy, a leading cause for emergent salpingectomy, affects approximately 1-2\% of all pregnancies globally, with rates remaining relatively stable over recent decades, though incidence can vary by region and risk factors. The incidence of salpingitis and hydrosalpinx, often sequelae of pelvic inflammatory disease (PID), has seen fluctuations influenced by trends in sexually transmitted infections, particularly among younger, sexually active women. Sterilization via bilateral salpingectomy or tubal ligation is a common contraceptive choice worldwide, with millions of procedures performed annually, making it one of the most frequent gynecologic surgeries. The demographic for sterilization typically includes women who have completed childbearing, often in their late 20s to 40s. The practice of opportunistic salpingectomy for ovarian cancer risk reduction is a more recent development, gaining significant traction in the last decade. While precise incidence rates are still being established, it is increasingly offered to women undergoing hysterectomy or other benign pelvic surgeries, particularly in Western countries, with uptake rates varying but steadily increasing. Mortality directly attributable to elective salpingectomy is exceedingly low, typically less than 0.1\%, reflecting the safety of modern surgical and anesthetic techniques. Morbidity, primarily related to general surgical risks like infection or bleeding, is also low, generally below 5\%, with serious complications being rare.

Outcomes depend on the indication and whether disease is unilateral or bilateral. Salpingectomy for ectopic pregnancy removes the affected tube but does not eliminate the risk of another ectopic pregnancy or guarantee fertility. Removal of a hydrosalpinx may improve the chance of IVF success in selected patients. Bilateral salpingectomy is intended to be permanent contraception, but exact failure risk is uncertain because long-term data are limited; it does not eliminate ovarian-cancer risk.

For cancer risk reduction, opportunistic salpingectomy is estimated to reduce the risk of high-grade serous ovarian cancer by a significant margin, potentially up to 50-65\% in the general population, and even higher in BRCA mutation carriers when combined with oophorectomy (salpingo-oophorectomy). This is supported by numerous observational studies and meta-analyses, leading to widespread adoption in clinical guidelines. Functional outcomes are generally very good, with most patients experiencing a full recovery and return to normal activities within weeks. Long-term prognosis is favorable, though potential for adhesion formation and chronic pelvic pain exists, albeit at low rates. Prognostic factors influencing outcomes include the underlying indication, the presence of comorbidities, the surgical approach (laparoscopic versus open), and the surgeon's experience. Landmark clinical trials and guidelines from organizations like ACOG and the Society of Gynecologic Oncology consistently support the efficacy and safety of salpingectomy for its various indications.

\section*{References}
\begin{enumerate}
  \item American College of Obstetricians and Gynecologists (ACOG). ACOG Practice Bulletin No. 191: Clinical Management of Ectopic Pregnancy. Obstet Gynecol. 2018;131(2):e65-e77.
  \item American College of Obstetricians and Gynecologists (ACOG). ACOG Practice Bulletin No. 193: Tubal Sterilization. Obstet Gynecol. 2018;131(4):e109-e120.
  \item Berek JS, Hacker NF. Berek \& Hacker's Gynecologic Oncology. 7th ed. Wolters Kluwer; 2021.
  \item MedlinePlus. Salpingectomy. U.S. National Library of Medicine; 2024. Available from: \url{https://medlineplus.gov/ency/article/002931.htm}
\end{enumerate}

\clearpage
